\documentclass[10pt,a4paper]{amsart}
\usepackage[margin=19mm]{geometry}
\usepackage{amsmath,amssymb,mathtools}
\usepackage[scr=rsfs,cal=euler]{mathalfa}
\usepackage{xfrac}
\usepackage[unicode,hidelinks,bookmarks=false]{hyperref}
\newcommand{\ie}{i.e.\ }
\newcommand{\cf}{cf.\ }
\newcommand{\resp}{respectively\ }
\newcommand{\Subsection}{\subsection}
\providecommand{\nobreakdash}{\nobreak\hbox{-}}
\setlength{\emergencystretch}{2em}
\multlinegap0pt
\usepackage{bm}
\usepackage{bbm}
\usepackage{dsfont}
\usepackage{xspace}
\usepackage{cancel}
\usepackage{mathtools}
\usepackage{amsthm}
\makeatletter
\@ifundefined{definition}{\newtheorem{definition}{Definition}}{}
\@ifundefined{lemma}{\newtheorem{lemma}[definition]{Lemma}}{}
\@ifundefined{theorem}{\newtheorem{theorem}[definition]{Theorem}}{}
\makeatother
\newcommand{\G}{\mathbb{G}}
\newcommand{\R}{\mathbb{R}}
\renewcommand{\phi}{\varphi}
\renewcommand{\vec}[1]{\bm{#1}}
\title[Additive relations in irrational powers]{Additive relations in irrational powers}
\author{Joseph Harrison}
\date{Source: arXiv:2512.04081v3 (CC BY 4.0)}
\begin{document}
\maketitle

\noindent\emph{Source and licence.} Additive relations in irrational powers. Version arXiv:2512.04081v3 (CC BY 4.0), distributed under the Creative Commons Attribution 4.0 International licence, as recorded in the arXiv licence metadata for this exact version and shown on the abstract page https://arxiv.org/abs/2512.04081v3. Full rights notice: licenses/arxiv-2512.04081-CC-BY-4.0.txt.\par\smallskip

\noindent\textit{Editorial scope.} Counted unit: the Lemma whose source label is \texttt{lemma/classification\_of\_semi-algebraic\_curves} in the article (source0.tex lines 683--693, source label \texttt{lemma/classification\_of\_semi-algebraic\_curves}), delivered together with the article's own complete original proof of it (source0.tex lines 694--725, one \texttt{proof} environment of 32 lines). No mathematics is written, repaired or re-derived: statement and proof are the article's own text line for line. Required context retained uncounted: functional\_transcendence\_theorem (lines 546--550). Every definition, display and label of those lines is retained unchanged. Omitted from this package and not redistributed: 1--545, 551--682, 726--985. Nothing inside a retained excerpt is omitted or altered: every retained body is delivered exactly as the article's lines are written, so no omission receipt of any kind accompanies this package. Citations: the retained excerpts carry no citation command at all, so no bibliography entry of the article is transcribed and no reference is redistributed. Reference adaptation: none was needed and none was made; every reference inside the delivered unit resolves inside it, so no unresolved reference or citation is produced. Printed slips kept literal: no printed word, symbol, hypothesis or number of the counted statement or of its proof was altered. Layout and class: the article's own class is \texttt{amsart} with packages including graphicx, amsmath, amssymb, amsfonts, mathrsfs, enumerate, lipsum, nomencl, tikz-cd, babel, xcolor, geometry, bm, hyperref; the wrapper is the lane's pinned \texttt{amsart} editorial wrapper at 10pt on A4, which restates the theorem-like environments the retained text needs and adds the article's own macro definitions that the retained text needs (\texttt{\textbackslash G}, \texttt{\textbackslash R}, \texttt{\textbackslash phi}, \texttt{\textbackslash vec}), with extra wrapper packages (bm, bbm, dsfont, xspace, cancel, mathtools). The delivered numbering is the unit's own. Assets: no retained span contains a figure environment, a graphics-inclusion command, a PDF-page inclusion, a table or a TikZ picture; no image file is copied into this package, the package declares no asset dependency and no image file is redistributed with it. Licence: version arXiv:2512.04081v3 of this article is distributed under the Creative Commons Attribution 4.0 International licence, as recorded in the arXiv licence metadata for this exact version and shown on the abstract page https://arxiv.org/abs/2512.04081v3. A full rights notice is delivered at licenses/arxiv-2512.04081-CC-BY-4.0.txt. Characters: every retained excerpt is pure ASCII, so the delivered engine, XeLaTeX with its default Latin Modern text font, reports no missing character.
\begin{theorem}
[Functional transcendence]
\label{functional_transcendence_theorem}
    Let $X$ be an irreducible, semi-algebraic subset of $\R_{>0}^n$. Let $\phi : \R_{>0}^n \to \R_{>0}^n$ be a continuous group homomorphism that does not restrict to a morphism of algebraic groups. Then the Zariski closure of $\phi(X)$ in $\G_m^n$ is a translate of a connected algebraic subgroup.
\end{theorem}

\begin{lemma}
\label{lemma/classification_of_semi-algebraic_curves}
    [Classification of semi-algebraic curves in $X_\lambda$]
    Let $\lambda = (c, \alpha, \vec{a}) \in \Lambda$ and suppose that $c$ is irrational. If $Z$ is an irreducible semi-algebraic curve in $X_\lambda$, then the Zariski closure of $Z$ has the form $t_{(\alpha, \dots, \alpha)}^*(\vec{g} H)$, where $H$ is a connected algebraic subgroup of $\G_m^{s + 1}$ of dimension $1$, $\vec{g} \in \R_{>0}^n$, and $t_{(\alpha, \dots, \alpha)}$ is the additive translation by the vector $(\alpha, \dots, \alpha)$. Moreover, there exists some $I \subseteq [s]$ with $1 \leq |I| \leq s - 1$ such that $H$ is defined by the equations $y_i = 1$ for $i \in I$ and $y_i = x$ for $i \in [s] \setminus I$, and such that the coordinates of the vector $\vec{g} = (g, g_1, \dots, g_s)$ satisfy the equations
    \begin{align*}
        \sum_{i \in I} g_i^c &= \sum_{i = 1}^{s - 1} a_i^c
        \\
        \sum_{i \in [s] \setminus I}
        g_i^c &= g^c.
    \end{align*}
\end{lemma}

\begin{proof}
    Let $Z$ be an irreducible semi-algebraic curve in $X_\lambda$, let $W = t_{(\alpha, \dots, \alpha)}(Z)$, and let
    \begin{align*}
        \phi(x, \vec{y}) = (x^c, y_1^c, \dots, y_s^c).
    \end{align*}
    Then, by Theorem \ref{functional_transcendence_theorem}, the Zariski closure of $\phi(W)$ is a translate $\vec{u} H$ of an connected algebraic subgroup $H$ of $\G_m^{s + 1}$ by some $\vec{u} \in \R_{>0}^{s + 1}$, lying in the linear subvariety
    \begin{align*}
        a_1^c + \dots + a_{s - 1}^c + x =
        y_1 + \dots + y_s.
    \end{align*}
    Let $\chi$ and $\chi_1, \dots, \chi_s$ denote the characters of $\G_m^{s + 1}$ corresponding to the projections to the coordinates $x, y_1, \dots, y_s$, respectively. If $\vec{u} = (u, u_1, \dots, u_s)$ then we have
    \begin{align*}
        a_1^c + \dots + a_{s - 1}^c + u\chi(h) =
        u_1\chi_1(h) + \dots + u_s\chi_s(h)
    \end{align*}
    for all $h \in H$. Suppose that $\chi$ is trivial on $H$. If each $\chi_i$ is also trivial on $H$ then $H$ is the trivial subgroup and $Z$ has dimension zero. Thus let $I \subseteq \{1, \dots, s\}$ be non-empty such that $i \in I$ if and only if $\chi_i$ is equal to a fixed, non-trivial character $\chi_I$ on $H$. By linear independence of characters we would then have
    \begin{align*}
        \sum_{i \in I} u_i = 0,
    \end{align*}
    which is absurd since $\vec{u} \in \R_{>0}^{s + 1}$. Therefore the character $\chi$ is non-trivial on $H$. By similar arguments, every character $\chi_1, \dots, \chi_s$ must be trivial or equal to $\chi$ on $H$. Let $I \subseteq [s]$ be such that $i \in I$ if and only if $\chi_i$ is trivial on $H$. Then $H$ is the $1$-dimensional subgroup defined by the equations $y_i = 1$ for all $i \in I$ and $y_i = x$ for all $i \in [s] \setminus I$, and linear independence of characters produces two relations
    \begin{align*}
        \sum_{i \in I} u_i &= \sum_{i = 1}^{s - 1} a_i^c, \\
        \sum_{i \in [s] \setminus I} 
        u_i &= u
    \end{align*}
    satisfied by the coordinates of $\vec{u}$. Now the Zariski closure of $\phi(W)$ is $\vec{u} H$ so $W$ is contained in $\phi^{-1}(\vec{u} H) = \phi^{-1}(\vec{u})H$. Let $\vec{g} = \phi^{-1}(\vec{u})$. Thus $Z$ is contained in $t_{(\alpha, \dots, \alpha)}^* (\vec{g}H)$ and if $\vec{g} = (g, g_1, \dots, g_s)$ then
    \begin{align*}
        \sum_{i \in I} g_i^c &= \sum_{i = 1}^{s - 1} a_i^c, \\
        \sum_{i \in [s] \setminus U} g_i^c &= g,
    \end{align*}
    as claimed. Since $Z$ and $t_{(\alpha, \dots, \alpha)}^* (\vec{g} H)$ are irreducible and of the same dimension, the Zariski closure of $Z$ must be as claimed. This completes the proof.
\end{proof}

\end{document}
